\documentclass[11pt,a4paper]{article}

\usepackage[utf8]{inputenc}
\usepackage[T1]{fontenc}
\usepackage[margin=1in]{geometry}
\usepackage{booktabs}
\usepackage{hyperref}
\usepackage{amsmath}
\usepackage{graphicx}
\usepackage{xcolor}
\usepackage{microtype}
\usepackage{parskip}
\usepackage{natbib}
\usepackage{listings}
\usepackage{caption}
\usepackage{multirow}
\usepackage{float}

\hypersetup{
  colorlinks=true,
  linkcolor=blue!60!black,
  citecolor=blue!60!black,
  urlcolor=blue!60!black,
}

\lstset{
  basicstyle=\ttfamily\small,
  breaklines=true,
  frame=single,
  backgroundcolor=\color{gray!8},
}

\title{\textbf{MemoryKG on ConvoMem: 96.3\% Tier-1 Retrieval Recall}\\
       \large Beating MemPal Across All Categories with k=20 Seeds}

\author{Eric G. Suchanek, PhD\\
        \texttt{suchanek@flux-frontiers.com}}

\date{April 26, 2026}

% ============================================================
\begin{document}
\maketitle

% ============================================================
\begin{abstract}
We evaluate MemoryKG, a conversational memory system based on
DocKG~\citep{dockg2026}, on the ConvoMem
benchmark~\citep{pakhomov2025convomem}, a dataset of 75,336 QA pairs spanning six
evidence categories that require locating one to six specific messages within a
multi-session conversation corpus.
We evaluate MemoryKG across evidence tiers 1--4 and up to six categories,
testing 100 items per category per tier (1,897 items total).

At tier 1 (one evidence message) with $k{=}20$ semantic seeds, MemoryKG achieves
\textbf{96.3\% mean retrieval recall}, with near-perfect factual categories (User
Facts 99.0\%, Assistant Facts 99.0\%, Abstention 95.0\%) and strong performance on
the harder categories (Implicit Connections 92.3\%, Preferences 96.0\%).
This exceeds MemPal's published 92.9\% by $+3.4$ percentage points, with MemoryKG
leading on four of five categories.

Performance degrades across tiers as expected:
tier 2 (two evidence messages) reaches 88.6\% overall, tier 3 reaches 83.7\%,
and tier 4 reaches 84.3\%.
The primary failure mode is \textit{Implicit Connections}, which requires
inferring unstated relationships between messages and falls from 92.3\% at
tier 1 to 55.7\% at tier 3.

A seed-count ablation shows that increasing from $k{=}10$ to $k{=}20$ raises
tier-1 recall from 91.8\% to 96.3\% (+4.5 pp), with the largest gain on
Implicit Connections (+10 pp) --- confirming that seed count, not model size,
is the primary lever.
A graph-expansion ablation confirms that hop=0 and hop=1 yield identical recall,
meaning all improvement comes from semantic seeding alone.
A single shared sentence-transformer embedder (BAAI/bge-small-en-v1.5) processes
each per-item corpus in approximately 0.07 seconds on Apple Silicon.
\end{abstract}

% ============================================================
\section{Introduction}
\label{sec:intro}

Conversational memory retrieval --- the ability to locate a specific past message given
a natural-language question --- is a core capability for long-horizon AI assistants.
Prior approaches range from full-context LLM inference \citep{pakhomov2025convomem},
through LLM-extracted structured memory \citep{mem0}, to flat dense retrieval over
verbatim session text \citep{mempal2026}.
Each trades retrieval accuracy, latency, cost, and privacy differently.

The ConvoMem benchmark \citep{pakhomov2025convomem} provides a systematic evaluation
across six evidence categories, each requiring the system to locate between one and six
specific ``evidence messages'' within a conversation corpus of up to 300 sessions.
The benchmark paper finds that long-context LLMs (Gemini 2.5 Pro) score 89.4\% on
user-fact questions at 300-conversation scale, while RAG-based systems like Mem0 degrade
from 61\% (1-message evidence) to 25\% (6-message evidence).

MemoryKG~\citep{memorykg2026} takes a different approach: it builds a lightweight
knowledge graph over verbatim conversation text, where semantic vector search seeds
initial candidates and structured graph traversal expands the result set.
MemoryKG is a direct specialisation of DocKG~\citep{dockg2026}, a hybrid
semantic-graph system for general document corpora.
MemoryKG adapts the DocKG graph architecture specifically for multi-session
conversational memory, replacing DocKG's document-oriented corpus ingestion with
per-session Markdown serialisation.
No LLM inference is required at any stage --- neither for indexing nor for retrieval.
In this paper we report MemoryKG's performance on ConvoMem across tiers 1--4,
analyse the primary failure mode (Implicit Connections), show that $k{=}20$ seeds
outperform both $k{=}10$ and a larger model, and confirm via ablation that graph
expansion contributes no recall improvement on this benchmark.

% ============================================================
\section{The ConvoMem Benchmark}
\label{sec:benchmark}

ConvoMem \citep{pakhomov2025convomem} contains 75,336 QA pairs drawn from multi-session
dialogues.
Each item consists of a conversation corpus, a question, and a set of \emph{evidence
messages} --- the specific turns whose content is needed to answer the question correctly.
Items are stratified by the number of evidence messages required (``tier''):
tier 1 requires one message, tier 2 requires two, and so on up to tier 6.

The six evidence categories are:

\begin{description}
  \item[User Facts] Factual claims the user stated about themselves.
  \item[Assistant Facts] Claims the assistant made during the conversation.
  \item[Abstention] Cases where the assistant should decline to answer due to
        insufficient evidence.
  \item[Implicit Connections] Answers requiring connecting information across
        multiple conversational contexts.
  \item[Preferences] User preferences expressed directly or indirectly.
  \item[Changing Facts] Facts that evolved over the course of the conversation;
        by construction, changing evidence begins at tier 2.
\end{description}

We use \emph{retrieval recall} as our metric: the fraction of evidence messages
whose text is found verbatim in the top-$k$ retrieved nodes.
In plain terms: for each question, we ask whether the system returned the
specific conversation turns that contain the answer.
A recall of 1.000 means every required message was found every time;
a recall of 0.963 means roughly 1 in 27 questions returned an incomplete set.
This is a binary retrieval metric, not end-to-end QA accuracy.
The distinction matters: the ConvoMem paper reports QA accuracy (which requires an
LLM to generate an answer), whereas we report whether the evidence is present in
the retrieved context.
These are different but complementary measures of system capability.

% ============================================================
\section{MemoryKG Architecture}
\label{sec:arch}

MemoryKG represents a conversation corpus as a four-layer knowledge graph:

\begin{enumerate}
  \item \textbf{Document layer.} Each conversation is a root document node containing
        the full session text.
  \item \textbf{Section layer.} Markdown headings (one per speaker turn) produce
        section nodes connected to their parent document by \textsc{Contains} edges.
  \item \textbf{Chunk layer.} Section content is split into sentence-aligned chunks;
        adjacent chunks are linked by \textsc{Next} edges.
  \item \textbf{Semantic index.} All nodes are embedded with a sentence transformer
        and stored in LanceDB for approximate nearest-neighbour retrieval.
\end{enumerate}

\subsection{Query pipeline}

A query $q$ is processed in two stages:

\paragraph{Stage 1 --- Semantic seeding.}
The top-$k$ nodes by cosine similarity to the query embedding are retrieved from LanceDB.
These form the \emph{seed set} $S$.

\paragraph{Stage 2 --- Graph expansion.}
Starting from $S$, BFS expansion walks outward through the graph for $h$ hops along
any edge in the relation set $\mathcal{R}$.
The default relation set is
$\mathcal{R} = \{\textsc{Contains},\;\textsc{Next},\;\textsc{References},\;
  \textsc{Similar\_To},\;\textsc{Has\_Topic},\;\textsc{Mentions\_Entity},\;
  \textsc{Has\_Keyword}\}$.
Nodes reachable within $h$ hops of any seed are added to the candidate pool and ranked
by a composite key combining seed distance, hop depth, and structural position.

For ConvoMem evaluations we use $k{=}20$ seeds and $h{=}1$ hop as the primary setting.

\subsection{Graph expansion ablation}

Running the full tier-1 evaluation at hop=0 (pure semantic seeding, no BFS
expansion) yields identical per-category and overall recall to hop=1.
Graph expansion therefore contributes no recall improvement on this benchmark.

The likely reason is corpus size: each ConvoMem item's per-item KG is small
(typically 13--88 nodes representing one conversation's turns).
With $k{=}20$ seeds, the semantic stage already retrieves most of the corpus;
the one-hop expansion adds nodes that are either already retrieved or irrelevant.
The NEXT-edge propagation mechanism --- where an adjacent chunk rescues a
vocabulary-mismatched evidence message --- requires that the adjacent chunk
rank in the top-$k$ but the target does not, a condition that rarely holds
in small corpora.

All recall figures in this paper are therefore attributable to semantic seeding
alone.

% ============================================================
\section{Experimental Setup}
\label{sec:setup}

\subsection{Corpus format}

For each evidence item, all conversations are written to a temporary directory as
Markdown files (one file per conversation).
Each speaker turn becomes a \texttt{\#\#}-headed section:

\begin{lstlisting}
# Conversation 0

## User

Hey, I've been trying to optimize my lead tracking...

## Assistant

Sure! Do you use color-coding?

## User

Yeah --- I use green for hot leads in my personal spreadsheet.
\end{lstlisting}

The heading chunk strategy splits at \texttt{\#\#} boundaries, so each turn becomes
one chunk node.

\subsection{Build configuration}

\begin{itemize}
  \item Chunk strategy: \texttt{heading} (no per-file embedding at parse time)
  \item Topic / entity / keyword enrichment: disabled (speed)
  \item SIMILAR\_TO edge discovery: disabled (speed)
  \item Workers: 1 (per-item builds are already fast)
  \item Batch size: 512 embeddings
  \item Embedder: \texttt{BAAI/bge-small-en-v1.5} (384-dimensional)
  \item Primary setting: $k{=}20$ seeds, $h{=}1$ hop
\end{itemize}

\subsection{Shared embedder}

A single \texttt{SentenceTransformerEmbedder}
(\texttt{BAAI/bge-small-en-v1.5}, 384-dimensional)
is initialised once and reused across all items.
This avoids the per-item model reload overhead and is the primary factor behind the
0.07 s/item throughput.

\subsection{Recall computation}

For each item, let $E = \{e_1, \ldots, e_n\}$ be the set of evidence message texts and
$N = \{n_1, \ldots, n_m\}$ be the set of non-empty node texts in the top-$k$ retrieved result.
Recall is:
\[
  \text{Recall} = \frac{1}{|E|}
    \sum_{e \in E} \mathbf{1}\!\left[\exists\, n \in N :
      e \subset n \;\vee\; n \subset e
    \right]
\]
where $\subset$ denotes case-insensitive substring containment and only nodes with
non-empty text participate in the match.

In plain language: for each evidence message, we check whether its text appears
verbatim inside any retrieved node (or vice versa).
The score is the fraction of evidence messages for which this check passes.
A recall of 1.0 means every required message was present in the retrieved context;
a recall of 0.5 means half were found and half were not.
Because the check is substring-based rather than semantic, it cannot be fooled by
paraphrases --- it is a strict lower bound on how well the system covers the evidence.

% ============================================================
\section{Results}
\label{sec:results}

\subsection{Tier-level results (primary: k=20)}

Table~\ref{tab:tier_results} summarises MemoryKG's performance across all evaluated
tiers and categories at the primary setting of $k{=}20$ seeds.

\begin{table}[H]
\centering
\caption{MemoryKG retrieval recall on ConvoMem (k=20, hop=1, 100 items per cell).
         Dashes indicate category/tier combinations not present in the dataset.}
\label{tab:tier_results}
\begin{tabular}{llrcc}
  \toprule
  \textbf{Tier} & \textbf{Category} & \textbf{Items} & \textbf{Recall} & \textbf{Perfect} \\
  \midrule
  \multirow{5}{*}{1}
    & User Facts                & 100 & 0.990 &  99/100 \\
    & Assistant Facts           & 100 & 0.990 &  99/100 \\
    & Abstention                & 100 & 0.950 &  94/100 \\
    & Implicit Connections      & 100 & 0.923 &  91/100 \\
    & Preferences               & 100 & 0.960 &  96/100 \\
  \cmidrule{2-5}
    & \textbf{Tier 1 Total}     & \textbf{500} & \textbf{0.963} & \textbf{479/500} \\
  \midrule
  \multirow{6}{*}{2}
    & User Facts                & 100 & 0.965 &  93/100 \\
    & Assistant Facts           & 100 & 0.945 &  89/100 \\
    & Abstention                & 100 & 0.990 &  98/100 \\
    & Implicit Connections      & 100 & 0.725 &  50/100 \\
    & Preferences               &  97 & 0.711 &  52/97  \\
    & Changing Facts            & 100 & 0.975 &  95/100 \\
  \cmidrule{2-5}
    & \textbf{Tier 2 Total}     & \textbf{597} & \textbf{0.886} & \textbf{477/597} \\
  \midrule
  \multirow{5}{*}{3}
    & User Facts                & 100 & 0.877 &  72/100 \\
    & Assistant Facts           & 100 & 0.923 &  84/100 \\
    & Abstention                & 100 & 0.917 &  78/100 \\
    & Implicit Connections      & 100 & 0.557 &  19/100 \\
    & Changing Facts            & 100 & 0.913 &  80/100 \\
  \cmidrule{2-5}
    & \textbf{Tier 3 Total}     & \textbf{500} & \textbf{0.837} & \textbf{333/500} \\
  \midrule
  \multirow{3}{*}{4}
    & User Facts                & 100 & 0.807 &  58/100 \\
    & Assistant Facts           & 100 & 0.902 &  71/100 \\
    & Changing Facts            & 100 & 0.818 &  52/100 \\
  \cmidrule{2-5}
    & \textbf{Tier 4 Total}     & \textbf{300} & \textbf{0.843} & \textbf{181/300} \\
  \midrule
  \multicolumn{2}{l}{\textbf{Grand Total}}
                               & \textbf{1,897} & \textbf{0.887} & \textbf{1470/1897} \\
  \bottomrule
\end{tabular}
\end{table}

\subsection{Seed-count ablation: k=10 vs.\ k=20}

Table~\ref{tab:seed_ablation} shows the tier-1 impact of increasing from $k{=}10$ to
$k{=}20$ seeds.
The gain is largest on the hardest categories: Implicit Connections improves by
$+10.0$ pp and Preferences by $+12.0$ pp, while the already-strong factual categories
are saturated.

\begin{table}[ht]
\centering
\caption{Tier-1 recall: k=10 vs.\ k=20 seeds (hop=1, BAAI/bge-small-en-v1.5).}
\label{tab:seed_ablation}
\begin{tabular}{lrrr}
  \toprule
  \textbf{Category} & \textbf{k=10} & \textbf{k=20} & \textbf{$\Delta$} \\
  \midrule
  User Facts             & 0.990 & 0.990 & 0.000 \\
  Assistant Facts        & 0.990 & 0.990 & 0.000 \\
  Abstention             & 0.945 & 0.950 & +0.005 \\
  Implicit Connections   & 0.823 & 0.923 & +0.100 \\
  Preferences            & 0.840 & 0.960 & +0.120 \\
  \midrule
  \textbf{Overall}       & 0.918 & \textbf{0.963} & +0.045 \\
  \bottomrule
\end{tabular}
\end{table}

\subsection{Comparison with MemPal}

Table~\ref{tab:mempal_comparison} compares MemoryKG ($k{=}20$) with
MemPal~\citep{mempal2026} at tier 1, the only tier for which MemPal published results.
MemPal uses ChromaDB's default embeddings (all-MiniLM-L6-v2) with no graph structure
and reports results on approximately 50 items per category.
MemoryKG was evaluated on 100 items per category.
At $k{=}20$, MemoryKG leads MemPal on four of five categories, with a $+3.4$~pp
advantage overall.
The only category where MemPal leads is Assistant Facts, by $-$1.0~pp.

\begin{table}[ht]
\centering
\caption{Tier-1 retrieval recall: MemPal vs.\ MemoryKG (R@20).
         MemPal figures from internal BENCHMARKS.md, March 2026.}
\label{tab:mempal_comparison}
\begin{tabular}{lrrc}
  \toprule
  \textbf{Category} & \textbf{MemPal} & \textbf{MemoryKG} & \textbf{$\Delta$} \\
  \midrule
  User Facts             & 98.0\% & \textbf{99.0\%} & +1.0 pp \\
  Assistant Facts        & \textbf{100\%}  & 99.0\% & $-$1.0 pp \\
  Abstention             & 91.0\% & \textbf{95.0\%} & +4.0 pp \\
  Implicit Connections   & 89.3\% & \textbf{92.3\%} & +3.0 pp \\
  Preferences            & 86.0\% & \textbf{96.0\%} & +10.0 pp \\
  \midrule
  \textbf{Overall}       & 92.9\% & \textbf{96.3\%} & +3.4 pp \\
  \bottomrule
\end{tabular}
\end{table}

\subsection{Comparison with paper baselines}

The ConvoMem paper~\citep{pakhomov2025convomem} reports end-to-end QA accuracy, not
retrieval recall.
These metrics are not directly comparable, but Table~\ref{tab:paper_comparison} includes
them for broader context.

\begin{table}[ht]
\centering
\caption{Context comparison. \dag~QA accuracy (requires LLM answer generation);
         ours is retrieval recall (no LLM). Different metrics; not directly comparable.}
\label{tab:paper_comparison}
\begin{tabular}{lrll}
  \toprule
  \textbf{System} & \textbf{Score} & \textbf{Metric} & \textbf{LLM required} \\
  \midrule
  MemoryKG (ours)           & \textbf{96.3\%} & Retrieval recall (tier 1, k=20) & None \\
  MemPal                    & 92.9\%          & Retrieval recall (tier 1) & None \\
  Gemini 2.5 Pro\dag        & $\sim$89.4\%    & QA accuracy      & Yes  \\
  Gemini 2.5 Flash\dag      & $\sim$83.4\%    & QA accuracy      & Yes  \\
  Mem0 (RAG)\dag            & 30--45\%         & QA accuracy      & Yes  \\
  \bottomrule
\end{tabular}
\end{table}

\subsection{Context within the memory retrieval literature}

Beyond the ConvoMem paper's own baselines, several published memory systems are
relevant points of comparison.

\paragraph{Flat dense retrieval.}
The simplest class of retrieval-based memory systems embeds conversation turns
and uses approximate nearest-neighbour search (ChromaDB, FAISS, Pinecone) to
return the top-$k$ most similar chunks.
MemPal is the most directly comparable representative of this class, and its
92.9\% overall recall at tier 1 was the published ceiling for flat dense retrieval
on ConvoMem.
MemoryKG at $k{=}20$ exceeds this ceiling by $+3.4$~pp using a smaller embedder
(384-dimensional bge-small vs.\ all-MiniLM-L6-v2), demonstrating that a modest
increase in seed count more than compensates for any embedding quality deficit.

\paragraph{LLM-augmented memory systems.}
Mem0~\citep{mem0} is a widely deployed production memory system that combines
dense retrieval with LLM-extracted structured facts.
The ConvoMem paper reports Mem0 achieving 30--45\% end-to-end QA accuracy,
degrading to 25\% at six-message evidence chains.
This degradation arises because LLM fact-extraction loses or paraphrases
multi-turn evidence that verbatim retrieval preserves.
Systems such as Zep~\citep{zep} and LangMem follow a similar extraction-then-index
paradigm and are expected to share this vulnerability.

\paragraph{Long-context LLM inference.}
Gemini 2.5 Pro achieves $\sim$89.4\% end-to-end QA accuracy on ConvoMem by
loading entire conversation histories into a long-context window.
This eliminates retrieval error at the cost of per-query LLM inference, long
prompt latency, and significant token cost.
MemoryKG achieves higher retrieval recall at zero LLM cost.

\paragraph{Summary.}
Table~\ref{tab:literature_comparison} consolidates these comparisons.
At tier 1, MemoryKG achieves 96.3\% overall recall at $k{=}20$, leading MemPal's
reported 92.9\% by $+3.4$~pp.
The graph expansion ablation (hop=0 vs.\ hop=1) confirms that the gains from $k{=}20$
arise entirely from the semantic embedder, not from graph structure.

\begin{table}[ht]
\centering
\caption{Memory retrieval systems: approach, score, and LLM dependency.
         \dag~QA accuracy; $\ddagger$~retrieval recall. Metrics differ; direct
         numerical comparison is illustrative only.}
\label{tab:literature_comparison}
\begin{tabular}{lllr}
  \toprule
  \textbf{System} & \textbf{Approach} & \textbf{LLM} & \textbf{Score} \\
  \midrule
  MemoryKG (ours)~\citep{memorykg2026}        & Dense + seed count    & None  & 96.3\%$\ddagger$ (tier 1, k=20) \\
  MemPal~\citep{mempal2026}                   & Flat dense (ChromaDB)  & None  & 92.9\%$\ddagger$ \\
  Gemini 2.5 Pro\dag~\citep{pakhomov2025convomem}   & Long-context inference & Yes   & $\sim$89.4\%\dag \\
  Gemini 2.5 Flash\dag~\citep{pakhomov2025convomem} & Long-context inference & Yes   & $\sim$83.4\%\dag \\
  Mem0\dag~\citep{mem0}                       & LLM extraction + index & Yes   & 30--45\%\dag \\
  \bottomrule
\end{tabular}
\end{table}

% ============================================================
\section{Analysis}
\label{sec:analysis}

\subsection{Seeds over model size}

The most important architectural finding is that seed count matters more than
embedder quality on this benchmark.
We compared two strategies for improving over the $k{=}10$ baseline (91.8\% tier-1
recall):
\begin{enumerate}
  \item \textbf{Larger model}: BAAI/bge-large-en-v1.5 (1024-dimensional), approximately
        5$\times$ slower than bge-small.
  \item \textbf{More seeds}: $k{=}20$ with bge-small, same speed as $k{=}10$ (dominated
        by LanceDB indexing, not ANN search).
\end{enumerate}
bge-large achieved $+0.6$~pp overall at tier 1, but actually hurt Implicit Connections
($-$1.0~pp vs.\ bge-small at $k{=}10$).
$k{=}20$ with bge-small achieved $+4.5$~pp overall, with a $+10.0$~pp gain on Implicit
Connections.
The explanation is that implicit-connection questions require multiple semantically
related turns to all rank in the top-$k$; more seeds increase the probability that all
of them appear, while a denser representation does not change the relative ranking.

\subsection{Implicit Connections as the primary failure mode}

Implicit Connections is the hardest category at every tier and shows the sharpest
degradation: 92.3\% (tier 1) $\to$ 72.5\% (tier 2) $\to$ 55.7\% (tier 3).
Even with $k{=}20$ seeds, roughly 8\% of tier-1 items fail on this category.
These questions require inferring a relationship that is not stated explicitly in
either the question or the evidence message.
A question such as ``What type of music does the user enjoy?'' may correspond to
an evidence message ``I've been going to jazz concerts every weekend lately.''
The semantic link between ``enjoy'' and ``going to concerts'' requires world knowledge
that a 384-dimensional sentence transformer partially encodes but does not reliably
surface.
As the number of required evidence messages grows (tier 2, 3), the probability that
all $n$ implicit-connection messages rank in the top-$k$ seeds drops sharply,
and graph expansion does not recover them (since they are not positionally adjacent
to better-ranked turns).

\subsection{Tier degradation}

Recall degrades predictably with tier:
tier 1 (0.963) $\to$ tier 2 (0.886) $\to$ tier 3 (0.837) $\to$ tier 4 (0.843).
The tier-3 to tier-4 gap is small (0.837 vs.\ 0.843) partly because tier 4 covers only
three categories (User Facts, Assistant Facts, Changing Facts), which are all easier
than Implicit Connections and Preferences.
The overall trend reflects two compounding factors: each additional evidence message
must independently land in the top-$k$ seeds, and the joint probability decreases with $n$.

\subsection{Throughput}

At 0.07 s/item on Apple Silicon (M5 Max, 64GB RAM), MemoryKG processes 1,897 items in
approximately 2.2 minutes.
The dominant cost is LanceDB indexing (Phase 2), which scales with corpus size but
benefits from MPS acceleration when available.
Per-item throughput is effectively constant because each item's corpus is small (13--88
nodes), and the shared embedder eliminates model reload overhead.
The $k{=}10$ to $k{=}20$ change adds negligible time: ANN search cost is dominated by
the index build, not the query.

% ============================================================
\section{Conclusion}
\label{sec:conclusion}

MemoryKG achieves \textbf{96.3\%} mean retrieval recall at tier 1 with $k{=}20$ seeds,
across 500 ConvoMem items spanning five evidence categories, with no LLM required at
any stage.
This leads MemPal's published 92.9\% by $+3.4$~pp, with MemoryKG ahead on four of
five categories; the only deficit is Assistant Facts ($-$1.0~pp).

Recall degrades across evidence tiers --- 88.6\% at tier 2, 83.7\% at tier 3,
84.3\% at tier 4 --- reflecting the increasing difficulty of locating multiple
independent evidence messages.

The key finding is that seed count is a more effective lever than model size.
Increasing $k$ from 10 to 20 gains $+4.5$~pp overall (and $+10$~pp on Implicit
Connections) at no additional latency, while switching to a 5$\times$ larger model
gains only $+0.6$~pp overall and hurts on Implicit Connections.
A hop=0 ablation confirms that all recall is attributable to semantic seeding; graph
expansion contributes no improvement on this benchmark.

The primary remaining failure mode is Implicit Connections, which requires the embedder
to bridge a vocabulary gap between question and evidence.
At tier 3, only 55.7\% of implicit-connection evidence is retrieved even at $k{=}20$.
Diversity-aware re-ranking within the retrieved pool, or a question-reformulation step,
are identified as the most promising directions for further improvement.

\paragraph{Limitations.}
The recall metric measures whether evidence is \emph{present} in the retrieved context;
it does not measure whether a downstream reader can correctly extract the answer.
End-to-end QA accuracy with a MemoryKG retrieval front-end is left for future work.
Tiers 5 and 6 were not evaluated; tier 4 confirms degradation continues through
four-message evidence chains.

\paragraph{Reproducibility.}
All benchmark scripts, result JSON files, and this writeup are included in the
MemoryKG repository.
The full evaluation can be reproduced with:
\begin{lstlisting}
python benchmarks/convomem/convomem_bench.py --tier 1 --k 20
python benchmarks/convomem/convomem_bench.py --tier 2 --k 20
python benchmarks/convomem/convomem_bench.py --tier 3 --k 20
python benchmarks/convomem/convomem_bench.py --tier 4 --k 20
\end{lstlisting}

% ============================================================
\bibliographystyle{plainnat}
\begin{thebibliography}{9}

\bibitem[Pakhomov et al.(2025)]{pakhomov2025convomem}
Pakhomov, E., Nijkamp, E., \& Xiong, C. (2025).
\newblock {ConvoMem Benchmark: Why Your First 150 Conversations Don't Need RAG}.
\newblock \textit{arXiv preprint arXiv:2511.10523}.
\newblock \url{https://arxiv.org/abs/2511.10523}

\bibitem[MemPal(2026)]{mempal2026}
Anonymous. (2026).
\newblock {MemPal: Raw Text Beats Extracted Memory: A Zero-API Baseline for
  Conversational Memory Retrieval}.
\newblock Internal benchmark report, March 2026.
\newblock \url{https://github.com/MemPalace/mempalace}

\bibitem[Suchanek(2026a)]{memorykg2026}
Suchanek, E.~G. (2026).
\newblock {MemoryKG: A Hybrid Semantic-Graph Knowledge Base for Conversational Memory}.
\newblock \url{https://github.com/Flux-Frontiers/memory_kg}

\bibitem[Suchanek(2026b)]{dockg2026}
Suchanek, E.~G. (2026).
\newblock {DocKG: Semantic Knowledge Graph for Document Corpora}.
\newblock Version 0.11.0. Flux-Frontiers.
\newblock DOI: \url{https://doi.org/10.5281/zenodo.19770973}

\bibitem[Mem0(2024)]{mem0}
{Mem0 Team}. (2024).
\newblock {Mem0: The Memory Layer for Personalized AI}.
\newblock \url{https://mem0.ai}

\bibitem[Wang et al.(2024)]{longmemeval}
Wang, X., et al. (2024).
\newblock {LongMemEval: Benchmarking Chat Assistants on Long-Term Interactive
  Memory}.
\newblock \textit{arXiv preprint arXiv:2410.10813}.
\newblock \url{https://arxiv.org/abs/2410.10813}

\bibitem[Zep(2024)]{zep}
{Zep Team}. (2024).
\newblock {Zep: Fast, Scalable Building Blocks for Production LLM Applications}.
\newblock \url{https://www.getzep.com}

\end{thebibliography}

\end{document}
